\chapter{Macro discrétionnaire}\label{ch:s2:discretionary-macro}

Un gérant macro qui anticipe une baisse des taux dispose d'une dizaine de façons de l'exprimer. La vue est la partie facile~; l'instrument, la taille et le stop
décident si avoir raison rapporte. Ce chapitre fixe une seule vue et fait varier tout le reste. Le gérant attend une baisse des taux de 50 points de base
sur six mois, et sur les trajectoires synthétiques le taux finit plus bas dans 73.4~\% des cas. Exprimée par des contrats à terme avec un stop à 25 points de base, la
vue est sortie sur 26.5~\% des trajectoires où elle se révèle juste. Avec un stop à 50 points de base, ce taux tombe à 5.6~\%. Exprimée par une option à la
monnaie, 17.2~\% des vues justes perdent tout de même la prime. La construction est \texttt{firm.macrobook}. Le chapitre ne comporte aucune fiche de stratégie~: son
sujet est la manière dont une vue macro devient une position.

\section{De la vue à la transaction}

\begin{definition}[Expression de la vue]\label{def:s2:discretionary-macro:expression}
Une \emph{expression de la vue}\index{expression de la vue} est l'instrument, le prix d'exercice, l'échéance et la taille choisis pour transformer une vue de marché en position~; des expressions différentes d'une même vue ne rapportent pas de la même façon,
selon l'ampleur, la vitesse et le chemin suivis par le marché.
\end{definition}

Une vue selon laquelle «~les taux vont baisser~» possède une taille (de combien), un horizon (d'ici quand), un chemin (régulier ou après une hausse) et un niveau de confiance. Chaque expression
parie sur une combinaison différente de ces éléments. Une position en contrats à terme parie sur la taille, quel que soit le chemin. Une option parie sur la taille et plafonne la perte à sa prime. Une
position en contrats à terme assortie d'un stop parie que le chemin ne part pas d'abord dans le mauvais sens. Un spread parie sur un mouvement jusqu'à un certain niveau et renonce au reste.

\texttt{firm.macrobook} rend ceci concret (le \cref{lst:s2:discretionary-macro:expressions}). La vue du gérant est une baisse de 50 points de base sur six
mois. La dérive réelle du taux est tirée autour de la vue avec une dispersion de 50 points de base, de sorte que la vue est juste en moyenne mais d'un montant variable.
Le taux évolue avec une volatilité normale de 90 points de base par an. Sur 20\,000 trajectoires, il finit plus bas dans 73.4~\% d'entre elles.

\section{Choix de l'instrument et convexité}

\begin{definition}[Budget de convexité]\label{def:s2:discretionary-macro:convexity}
Un \emph{budget de convexité}\index{budget de convexité} est la prime qu'un gérant accepte de dépenser en options pour exprimer ses vues, dont il sait que c'est la perte maximale~;
dépenser ce budget permet d'obtenir des positions dont la perte est connue à l'avance et dont le gain croît avec l'ampleur du mouvement.
\end{definition}

Chaque expression est dimensionnée sur le même budget de risque. Pour les options, le budget est la prime. Pour les contrats à terme assortis d'un stop, le budget est la distance au stop. Pour les contrats à terme
sans stop, il s'agit d'une perte à deux écarts types sur l'horizon. Les options sont évaluées avec le modèle normal, à la volatilité du taux augmentée de 10~\%. Une
option receveuse à la monnaie coûte 27.9 points de base, une option dont le prix d'exercice est inférieur de 25 points de base coûte 17.2, et un spread d'options entre les prix d'exercice 0 et
$-50$ coûte 18.2 (la \cref{fig:s2:discretionary-macro:expressions}).

\begin{center}
\small
\setlength{\tabcolsep}{4pt}%
\begin{tabular}{@{}p{4.4cm}ccccc@{}}
\toprule
20\,000 trajectoires, par unité de budget de risque & moyenne & écart-type & P(gain) & P(perte $|$ juste) & 5\up{e}~centile\\
\midrule
contrats à terme, sans stop & 0.39 & 0.63 & 73~\% & 0.0~\% & $-0.65$\\
contrats à terme, stop 25 pb & 1.68 & 2.98 & 54~\% & 26.5~\% & $-1.00$\\
contrats à terme, stop 50 pb & 0.98 & 1.55 & 69~\% & 5.6~\% & $-1.00$\\
option à la monnaie & 1.26 & 2.23 & 61~\% & 17.2~\% & $-1.00$\\
option hors de la monnaie de 25 pb & 1.69 & 3.20 & 54~\% & 26.5~\% & $-1.00$\\
spread d'options (0 à $-50$) & 0.71 & 1.22 & 65~\% & 10.8~\% & $-1.00$\\
\bottomrule
\end{tabular}
\end{center}

Le classement dépend de ce que représente le budget. Par unité de perte à deux écarts types, les contrats à terme rapportent 0.39 et ne perdent jamais lorsque la vue est juste. Par unité
de prime, l'option hors de la monnaie rapporte 1.69. Ce chiffre tient surtout à l'effet de levier, avec la même part de 26.5~\% de vues justes qui expirent sans valeur, comme pour le
stop serré. Le spread d'options est l'option la plus prudente~: il renonce aux grands mouvements, que la vue attend moins, et gagne plus souvent (65~\%). Aucune
expression n'est meilleure qu'une autre en général. Chacune convient à une conviction différente quant à la manière dont la vue se réalisera.

\begin{omfigure}
\begin{tikzpicture}
\begin{axis}[width=12cm,height=5.6cm,ybar,bar width=11pt,ylabel={\small par unité de budget / part},tick label style={font=\footnotesize},
  xmin=0.4,xmax=6.6,ymin=0,ymax=1.9,xtick={1,2,3,4,5,6},
  xticklabels={contrats à terme,{stop 25 pb},{stop 50 pb},{option ATM},{option OTM},{spread}},x tick label style={font=\scriptsize,rotate=20,anchor=north east},
  grid=major,grid style={gray!25},table/col sep=comma,legend style={font=\scriptsize,at={(0.5,-0.36)},anchor=north},legend columns=2]
\addplot[fill=omThm!70,draw=omThm] table[x=x,y=mean]{figdata/strategies-2/17-discretionary-macro/expressions.csv};
\addplot[fill=gray!40,draw=gray] table[x=x,y=lose]{figdata/strategies-2/17-discretionary-macro/expressions.csv};
\legend{gain moyen par unité de budget de risque,part des vues justes qui perdent}
\end{axis}
\end{tikzpicture}

\omcaption{Une même vue (une baisse de 50 points de base en six mois) exprimée de six façons sur 20\,000 trajectoires synthétiques~: le gain moyen par unité du budget de risque de chaque
expression, et la part des trajectoires où la vue était juste mais où l'expression a perdu. Données~: \texttt{s2\_macrobook.table}.}\label{fig:s2:discretionary-macro:expressions}
\end{omfigure}

\section{Dimensionnement et stops}

\begin{definition}[Discipline des stops]\label{def:s2:discretionary-macro:stop}
La \emph{discipline des stops}\index{discipline des stops} est la règle, fixée à l'ouverture d'une position, qui prescrit de la clôturer après une perte ou un mouvement défavorable donné, ainsi que
la pratique consistant à suivre cette règle~; elle limite la perte sur les vues fausses au prix de la clôture de certaines vues justes avant qu'elles ne rapportent.
\end{definition}

Un stop transforme le chemin en risque. Le taux peut monter de 25 points de base avant de baisser de 50, et un stop à 25 clôt la position avant la baisse. La part des vues justes sorties par le stop
diminue fortement avec la distance du stop (la \cref{fig:s2:discretionary-macro:stops})~: 55.6~\% à 10 points de base,
26.5~\% à 25, 5.6~\% à 50, 0.7~\% à 75 et 0.1~\% à 100. Kaminski et Lo ont proposé un cadre permettant de déterminer quand les règles de stop-loss ajoutent ou retranchent du rendement espéré
et de la volatilité. Dans leur analyse empirique des contrats à terme sur indices, certaines politiques de stop-loss ont augmenté le rendement espéré tout en réduisant la volatilité à des fréquences
d'échantillonnage plus longues. Sur le marché synthétique, la dérive est la vue elle-même, de sorte que le stop ne rapporte que lorsque la vue est fausse.

\begin{omfigure}
\begin{tikzpicture}
\begin{axis}[width=10cm,height=5.4cm,xlabel={\small distance du stop (pb)},ylabel={\small vues justes sorties par le stop (\%)},tick label style={font=\footnotesize},
  xmin=0,xmax=100,ymin=0,ymax=75,grid=major,grid style={gray!25},table/col sep=comma]
\addplot[omThm,thick,mark=*,mark size=1.3pt] table[x=stop,y=share]{figdata/strategies-2/17-discretionary-macro/stops.csv};
\end{axis}
\end{tikzpicture}

\omcaption{La part des vues justes (le taux finit plus bas après six mois) dont la position en contrats à terme a été sortie en premier par le stop, selon la distance du stop,
sur les trajectoires synthétiques. Données~: \texttt{s2\_macrobook.stops}.}\label{fig:s2:discretionary-macro:stops}
\end{omfigure}

Le dimensionnement découle du budget et du stop. À budget égal, un stop plus serré signifie une position plus grande~: le stop à 25 points de base porte
deux fois la position du stop à 50. C'est pourquoi son gain moyen par unité de budget est plus élevé (1.68 contre 0.98), alors même qu'il est sorti beaucoup plus souvent. Un gérant
dont le budget est une limite de perte mensuelle (livre~8, chapitre~28, sur les limites de drawdown) doit choisir entre une petite position avec un stop large et une
grande position avec un stop serré. La première résiste au bruit~; la seconde a raison moins souvent, mais rapporte davantage lorsqu'elle a raison.

\section{Tenir le score}

Un carnet discrétionnaire se juge sur ses résultats, alors que ses décisions se prennent à partir de vues. Séparer les deux est précisément l'intérêt de tenir un score. Il faut consigner chaque vue
au moment où elle se forme~: taille, horizon, niveau de confiance. Il faut noter l'expression choisie et la raison de ce choix. À la clôture de la transaction, il faut noter si la vue s'est réalisée,
si l'expression a rapporté, et si c'est le stop ou l'horizon qui a mis fin à la position.

Sur un grand nombre de transactions, le registre montre si les vues du gérant valent mieux que le hasard, et si les expressions tirent le meilleur parti de ces vues. Sur le
marché synthétique, avec une vue connue exactement (une dispersion nulle), le taux finit plus bas dans 78.6~\% des trajectoires~: le reste n'est que du bruit sur six mois.
Même une vue parfaitement informée n'a raison que dans cette proportion à cet horizon. Les contrats à terme assortis d'un stop serré rapportent alors 1.45 par unité de budget, et l'option à la
monnaie 1.07. Un score fondé sur la seule distinction entre juste et faux ne permettrait pas de le voir.

\section{Tutoriel~: avoir raison et perdre quand même}\label{tut:s2:discretionary-macro:tut}

\begin{tutorial}
\textbf{Objectif.} Simuler les trajectoires d'une vue, l'exprimer de six façons, et mesurer à quelle fréquence une vue juste perd et de combien les stops la raccourcissent.
\textbf{État final~:} le tableau et les deux figures.
\begin{tutsteps}
\item \textbf{Expressions}.
\omcode{code/firm/macrobook/firm_macrobook.py}{60}{82}{Gain de chaque expression par unité de budget de risque, et vues justes sorties par stop.}\label{lst:s2:discretionary-macro:expressions}
\item \textbf{Le tableau}.
\omcode{code/strategies-2/17-discretionary-macro/python/s2_macrobook.py}{30}{39}{Moyenne, fourchette, probabilité de gain et de perte quand la vue est juste.}\label{lst:s2:discretionary-macro:table}
\item \textbf{Exécuter} \texttt{table()}, \texttt{stops()}, \texttt{premiums()} et \texttt{fig\_macrobook.py}.
\end{tutsteps}
\textbf{Pour aller plus loin.} Rendre incertain le moment de la baisse (elle survient dans la seconde moitié)~; ajouter un stop suiveur~; exprimer la vue comme
un trade de courbe dans un modèle à deux facteurs.
\end{tutorial}

\section{Construction~: livre macro}\label{bld:s2:discretionary-macro:macrobook}

\begin{build}
\bfield{Rôle} Une distribution des trajectoires d'une vue, des expressions dimensionnées sur un budget de risque commun, et des règles de stop avec leur coût.
\bfield{Interface} \texttt{ViewConfig(\dots)}, \texttt{rate\_paths(cfg)}, \texttt{bachelier\_receiver(strike, vol, tau)},
\texttt{expressions(paths, cfg, stops)}, \texttt{stopped\_correct(paths, cfg, stop)}.
\bfield{Règles} Options évaluées à la volatilité de la trajectoire multipliée par une prime~; gain de chaque expression divisé par son propre budget de risque.
\bfield{Tests d'acceptation} \texttt{code/firm/macrobook/tests/}~: le prix du modèle normal à la monnaie et la parité call-put~; expressions et stops sur
des trajectoires construites à la main.
\bfield{Pour aller plus loin} Calendrier incertain~; stops suiveurs~; courbes à deux facteurs.
\end{build}

\begin{omsources}
\item K.~M.~Kaminski et A.~W.~Lo, ``When do stop-loss rules stop losses?'', \emph{Journal of Financial Markets} 18, 2014.
\end{omsources}

\section{Exercices}

\begin{exercise}[$\star$]\label{exo:s2:discretionary-macro:1}
Calculer le prix d'une option de vente à la monnaie sur un taux dont la volatilité normale est de 99 points de base par an, sachant qu'il reste six mois avant l'échéance.
\end{exercise}

\begin{exercise}[$\star$]\label{exo:s2:discretionary-macro:2}
Un budget d'une unité achète des contrats à terme avec un stop à 25 points de base ou avec un stop à 50 points de base. Comparer la taille des deux positions.
\end{exercise}

\begin{exercise}[$\star$]\label{exo:s2:discretionary-macro:3}
Le taux baisse de 80 points de base. Calculer ce que rapporte le spread d'options (0 à $-50$) par unité de sa prime de 18.2 points de base.
\end{exercise}

\begin{exercise}[$\star\star$]\label{exo:s2:discretionary-macro:4}
Expliquer pourquoi la part des vues justes sorties par stop est si sensible à la distance du stop.
\end{exercise}

\begin{exercise}[$\star\star$]\label{exo:s2:discretionary-macro:5}
Expliquer pourquoi l'option hors de la monnaie offre un gain moyen par unité de budget plus élevé, et une probabilité de gain plus faible, que l'option à la monnaie.
\end{exercise}

\begin{exercise}[$\star\star$]\label{exo:s2:discretionary-macro:6}
Expliquer pourquoi une vue exactement juste reste fausse sur 21~\% des trajectoires à six mois.
\end{exercise}

\begin{exercise}[$\star\star\star$]\label{exo:s2:discretionary-macro:7}
\emph{Programmation.} Exécuter \texttt{table(0.0)}, une vue sans dispersion. Comment les probabilités et les gains changent-ils, et qu'est-ce que cela dit de la
valeur de la précision~?
\end{exercise}

\begin{exercise}[$\star\star\star$]\label{exo:s2:discretionary-macro:8}
\emph{Trouver la faille.} «~J'ai eu raison sur les taux cette année et j'ai quand même perdu de l'argent, donc mes vues sont bonnes et le marché a eu tort.~»
\end{exercise}

\section{Problème~: avoir raison et perdre quand même}

\begin{problem}[{Problème du week-end --- exprimer une vue macro}]\label{pb:s2:discretionary-macro:1}
Les trajectoires synthétiques du chapitre et les données publiques.

\medskip\noindent\textbf{Partie I --- La vue.}
\begin{enumerate}
\item Définir une \omterm{def:s2:discretionary-macro:expression}{expression de la vue}.
\item Sur quoi chaque expression parie-t-elle~?
\item Décrire la vue synthétique et ses trajectoires.
\item À quelle fréquence la vue est-elle juste~?
\end{enumerate}

\medskip\noindent\textbf{Partie II --- Instruments.}
\begin{enumerate}[resume]
\item Définir un \omterm{def:s2:discretionary-macro:convexity}{budget de convexité}.
\item Comment chaque expression est-elle dimensionnée~?
\item Donner les prix des options.
\item Pourquoi le classement dépend-il du budget~?
\end{enumerate}

\medskip\noindent\textbf{Partie III --- Stops.}
\begin{enumerate}[resume]
\item Définir la \omterm{def:s2:discretionary-macro:stop}{discipline des stops}.
\item Donner la part des vues justes sorties par stop selon la distance du stop.
\item Qu'ont étudié et montré Kaminski et Lo~?
\item Comment le stop et la taille interagissent-ils~?
\end{enumerate}

\medskip\noindent\textbf{Partie IV --- Le verdict.}
\begin{enumerate}[resume]
\item Énoncer le \emph{résultat nommé}~: le gain espéré de chaque expression d'une même vue et la part des vues justes sorties tôt par stop.
\item Quelle expression convient à une vue sur la taille, et laquelle à une vue sur la trajectoire~?
\item Qu'est-ce qui change lorsque la vue est exacte~?
\item Que doit noter un gérant pour chaque transaction~?
\item Comment juger un gérant discrétionnaire~?
\item Pourquoi ce chapitre ne contient-il aucune fiche de stratégie~?
\item Quel est le lien entre ce chapitre et le chapitre~16~?
\item En une phrase~: qu'est-ce qui décide si une vue juste rapporte~?
\end{enumerate}
\end{problem}

\section{Questions d'entretien}

\begin{interviewq}[$\star$ \iqroles{trader}]\label{iq:s2:discretionary-macro:1}
Vous pensez que la Fed baissera ses taux dans six mois. Comment l'exprimeriez-vous~?
\end{interviewq}

\begin{interviewq}[$\star\star$ \iqroles{trader}]\label{iq:s2:discretionary-macro:2}
Quand utiliseriez-vous des options plutôt que des contrats à terme pour une vue macro~?
\end{interviewq}

\begin{interviewq}[$\star\star$ \iqroles{risque}]\label{iq:s2:discretionary-macro:3}
Comment fixeriez-vous les stops d'un livre macro discrétionnaire~?
\end{interviewq}

\begin{interviewq}[$\star\star$ \iqroles{chercheur}]\label{iq:s2:discretionary-macro:4}
Comment mesureriez-vous si les vues d'un gérant ont du talent~?
\end{interviewq}

\begin{interviewq}[$\star\star$ \iqroles{développeur}]\label{iq:s2:discretionary-macro:5}
Concevez un journal de transactions qui consigne les vues, les expressions et les résultats.
\end{interviewq}

\begin{interviewq}[$\star\star\star$ \iqroles{chercheur}]\label{iq:s2:discretionary-macro:6}
Pour un mouvement brownien sans dérive de volatilité $\sigma$, établir la probabilité d'atteindre $+s$ avant l'instant $T$, et s'en servir pour expliquer la forme de la
courbe de sorties par stop du chapitre.
\end{interviewq}
